
\chapter{Part IV Review: From Modes to Particles}

The free scalar field is a continuum limit of many coupled oscillators.  Quantization promotes the mode amplitudes to operators, and the creation operators build Fock space.

\section{A guided reconstruction}

% QFTFIGURE-BEGIN partreview04-01-part-map
\qftfigure{figures/instructional/partreview04/partreview04-01-part-map.pdf}{The map collects the part's main constructions into a visual route so the reader can see how the chapters support one another.}{fig:partreview04-01-part-map}
% QFTFIGURE-END partreview04-01-part-map












Start in a finite box, where momenta are discrete.  Normalize the modes so that
\[
  [a_{\bm p},a_{\bm q}^\dagger]=\delta_{\bm p\bm q}.
\]
Only after this finite normalization is clear should one take the continuum limit and replace Kronecker deltas by Dirac deltas.  The propagator is the Green function of the quantized field equation with the Feynman boundary condition.

\begin{exercise}
Why is it safer to quantize in a box before writing continuum delta functions?
\begin{answer}
The box gives ordinary sums and Kronecker deltas, making normalization explicit before the continuum limit introduces distributions.
\end{answer}
\end{exercise}
